\section{Metodologia}

\begin{frame}{Implementação e ferramentas}
  \begin{columns}[T,totalwidth=\textwidth]
    \begin{column}{0.55\textwidth}
      \begin{itemize}
        \item \textbf{Linguagem/framework:} Python + \textbf{PyTorch} (diferenciação
              automática, GPU);
        \item \textbf{Operadores neurais} (FNO/PINO): adaptados da biblioteca
              \textbf{\emph{neuraloperator}} --- camada \texttt{SpectralConv2d} com FFT e
              truncamento de modos;
        \item \textbf{Solver de referência:} implementação pseudoespectral própria
              (FFT + RK4) que gera o \emph{dataset} e serve de \emph{baseline};
        \item \textbf{PINN:} MLP com $\tanh$, resíduos por diferenciação automática.
      \end{itemize}
    \end{column}
    \begin{column}{0.45\textwidth}
      \begin{information}[Hardware modesto]
        Todos os experimentos rodaram numa única GPU \textbf{NVIDIA GTX 1660 SUPER} --- uma
        placa de entrada. Um objetivo prático é mostrar que esses métodos são viáveis
        \alert{sem infraestrutura de ponta}.
      \end{information}
      \begin{remark}
        Reprodutibilidade garantida por \emph{seed} fixa ($37$) e modo determinístico do
        cuDNN.
      \end{remark}
    \end{column}
  \end{columns}
\end{frame}

\begin{frame}{Design dos experimentos}
  Os experimentos são desenhados para responder às três perguntas dos objetivos:
  \begin{columns}[T,totalwidth=\textwidth]
    \begin{column}{0.5\textwidth}
      \begin{block}{Variantes (dados vs.\ física)}
        \footnotesize
        \begin{itemize}
          \item \textbf{FNO}: puramente \emph{data-driven};
          \item \textbf{PINO}: com dados \emph{e} sem dados;
          \item \textbf{PINN}: com dados \emph{e} sem dados.
        \end{itemize}
      \end{block}
      \begin{block}{Regimes (acurácia vs.\ dificuldade)}
        \footnotesize
        Família $\alpha=\beta$ com $20$ amostras em $[0{,}1,\,4{,}0]$; avaliação em
        $\alpha=\beta \in \{0{,}1,\ 3{,}21,\ 4{,}2\}$ (inclui \emph{extrapolação}).
      \end{block}
    \end{column}
    \begin{column}{0.5\textwidth}
      \begin{block}{Resolução (independência de malha)}
        \footnotesize
        Treino em $N_x=128$; avaliação também em $256$ e $512$ para testar
        \emph{super-resolution}.
      \end{block}
      \begin{block}{Métricas}
        \footnotesize
        \begin{itemize}
          \item erro relativo $L^2$ no espaço-tempo;
          \item erro relativo no \alert{espectro de Fourier} (diagnóstico do viés espectral);
          \item evolução do erro por frequência ao longo do treino.
        \end{itemize}
      \end{block}
    \end{column}
  \end{columns}
\end{frame}

\begin{frame}{Configurações e tamanho dos modelos}
  \begin{columns}[c,totalwidth=\textwidth]
    \begin{column}{0.55\textwidth}
      \footnotesize
      \setlength{\tabcolsep}{4pt}
      \centering
      \begin{tabular}{@{}lcc@{}}
        \toprule
        & \textbf{FNO / PINO} & \textbf{PINN} \\
        \midrule
        Épocas              & $15{,}000$    & $15{,}000$ \\
        Otimizador          & Adam          & Adam \\
        \emph{Learning rate}& $10^{-3}$     & $10^{-3}$ \\
        Modos $(k_x,k_t)$   & $16\times 16$ & --- \\
        Largura/neur.       & $32$          & $256$ \\
        Camadas             & $4$           & $4$ \\
        Pontos/\emph{batch} & $256$         & $5000$+$500$ \\
        \bottomrule
      \end{tabular}
    \end{column}
    \begin{column}{0.45\textwidth}
      \begin{block}{Tamanho dos modelos}
        \footnotesize
        \begin{itemize}
          \item \textbf{PINN}: $\approx 300$K parâmetros;
          \item \textbf{FNO / PINO}: $\approx 1$M parâmetros.
        \end{itemize}
      \end{block}
      \begin{block}{Pesos da perda (PINO)}
        \footnotesize
        $\lambda_{\text{pde}}=100,\ \lambda_{\text{ic}}=1,\ \lambda_{\text{data}}=1$.
      \end{block}
      {\footnotesize Domínio: $x\in[-60,60]$, $t\in[0,30]$.}
    \end{column}
  \end{columns}
\end{frame}
